\documentclass[../../main.tex]{subfiles}
\begin{document}


{\bf [解]}

\begin{figure}[htb]
\begin{center}
 % --- 1. 0 < a <= 4/3 の時 (例: a=1) ---
\begin{tikzpicture}[scale=0.5, >=stealth, every node/.style={font=\scriptsize}]
  % 解析パラメータ (a=1) -> A1(6,0), A2(8,0), B1(0,4), B2(0,3.2)
  % D_a: OA1PB2
  
  % 塗りつぶし領域 D_a
  \fill[pattern=north east lines, pattern color=gray!60] (0,0) -- (6,0) -- (3,2) -- (0,3.2) -- cycle;

  % 座標軸
  \draw[->] (-0.5,0) -- (8.5,0) node[right] {$x$};
  \draw[->] (0,-0.5) -- (0,5.5) node[above] {$y$};
  \node[below left] at (0,0) {$O$};

  % 直線 ell_1: 2x+3y=12 -> 切片 (6,0), (0,4)
  \draw[thick] (6,0) -- (0,4);
  \node[blue, anchor=south west] at (1,3.3) {$\ell_1$};
  
  % 直線 ell_2: x + 2.5y = 8 (a=1) -> 切片 (8,0), (0,3.2)
  \draw[thick, red] (8,0) -- (0,3.2);
  \node[red, anchor=north east] at (7.5,0.8) {$\ell_2$};

  % 定点 P(3,2)
  \fill (3,2) circle (2.5pt);
  \node[above right] at (3,2) {$P(3,2)$};

  % 頂点ラベル
  \node[below] at (6,0) {$A_1$};
  \node[below, red] at (8,0) {$A_2$};
  \node[left] at (0,4) {$B_1$};
  \node[left, red] at (0,3.2) {$B_2$};

\end{tikzpicture}
\caption{二直線$\ell_1,\ell_2$と領域$D_a$（斜線部）の様子}
\end{center}
\end{figure}


$a>0$のもとで，
\begin{align}
 \begin{cases}
  x\ge 0 \\
  y\ge 0 \\
  2x+3y\le 12 \\
  ax+\left(4-\dfrac{3a}{2}\right)y\le 8
 \end{cases} \label{eq:1}
\end{align}
の表す領域を$D_a$とし，$D_a$上で$g(x,y)=x+y$の最大値$f(a)$をもとめる．\cref{eq:1}の後者二つの領域の境界の直線を$l_1,l_2$とおく：
\begin{align}
 \begin{cases}
  \ell_1: 2x + 3y - 12 = 0 \\
  \ell_2: ax + \left(4-\dfrac{3a}{2}\right)y -8 = 0
 \end{cases} \label{eq:2}
\end{align}
と，これらは$a$の値によらず定点$P(3,2)$を通る．よって$P$は常に$D_a$上にあり，
\begin{align}
 g(x,y) = g(3,2) = 5 \label{eq:3}
\end{align}
である．
両者は直線で，かつ$a$を動かしても$\ell_2$が$\ell_1$と一致することはないから交点はこの点のみである．
従って，$D_a$は原点$O(0,0)$，$x$軸上の$\ell_1,\ell_2$の切片$A_1,A_2$のうち$x$座標の小さい方，$P$，$y$軸上の$\ell_1,\ell_2$の切片$B_1,B_2$のうち$y$座標の小さい方の4点からなる四角形をなす．
\cref{eq:2}より
\begin{align}
 A_1(6,0), A_2\left(\dfrac{8}{a},0\right), B_1\left(0,4\right), B_2\left(0,\dfrac{8}{4-\frac{3a}{2}}\right) \label{eq:4}
\end{align}
である．ただし$B_2$は$a=\dfrac{8}{3}$の時には存在しない．
従って$a$の値に応じて以下のように領域が切り替わる．
$g$も線形だからその最大値は$D_a$の端点でとることに注意する．


\subparagraph{$1^\circ$ $0<a\le\dfrac{4}{3}$の時}

\begin{figure}[htb]
\begin{center}
 % --- 1. 0 < a <= 4/3 の時 (例: a=1) ---
\begin{tikzpicture}[scale=0.5, >=stealth, every node/.style={font=\scriptsize}]
  % 解析パラメータ (a=1) -> A1(6,0), A2(8,0), B1(0,4), B2(0,3.2)
  % D_a: OA1PB2
  
  % 塗りつぶし領域 D_a
  \fill[pattern=north east lines, pattern color=gray!60] (0,0) -- (6,0) -- (3,2) -- (0,3.2) -- cycle;

  % 座標軸
  \draw[->] (-0.5,0) -- (8.5,0) node[right] {$x$};
  \draw[->] (0,-0.5) -- (0,5.5) node[above] {$y$};
  \node[below left] at (0,0) {$O$};

  % 直線 ell_1: 2x+3y=12 -> 切片 (6,0), (0,4)
  \draw[thick] (6,0) -- (0,4);
  \node[blue, anchor=south west] at (1,3.3) {$\ell_1$};
  
  % 直線 ell_2: x + 2.5y = 8 (a=1) -> 切片 (8,0), (0,3.2)
  \draw[thick, red] (8,0) -- (0,3.2);
  \node[red, anchor=north east] at (7.5,0.8) {$\ell_2$};

  % 定点 P(3,2)
  \fill (3,2) circle (2.5pt);
  \node[above right] at (3,2) {$P(3,2)$};

  % 頂点ラベル
  \node[below] at (6,0) {$A_1$};
  \node[below, red] at (8,0) {$A_2$};
  \node[left] at (0,4) {$B_1$};
  \node[left, red] at (0,3.2) {$B_2$};

\end{tikzpicture}
\caption{$0<a\le\dfrac{4}{3}$の時の$D_a$．}
\end{center}
\end{figure}

この時は$6 \le \dfrac{8}{a}$だから$A_2$よりも$A_1$が原点に近く，また$\dfrac{8}{4-\frac{3a}{2}} \le 4$より$B_2$が$B_1$より原点に近いので$D_a$は四角形$OA_1PB_2$になる．
頂点での$g(x,y)$の値は
\begin{align}
 \begin{cases}
  g(A_1) = g\left(6,0\right) = 6 \\
  g(P)   = g(3,2) = 5 \\
  g(B_2) = g\left(0,\dfrac{8}{4-\frac{3a}{2}}\right) = \dfrac{8}{4-\frac{3a}{2}}
 \end{cases}
\end{align}
となるから，最大値は
\begin{align}
 \max g(x,y) = g(A_1) = 6 \label{eq:5}
\end{align}
である．

\subparagraph{$2^\circ$ $\dfrac{4}{3}\le a < \dfrac{8}{3}$の時}

\begin{figure}[htb]
 \begin{center}
  \begin{tikzpicture}[scale=0.5, >=stealth, every node/.style={font=\scriptsize}]
  % 解析パラメータ (a=2) -> A1(6,0), A2(4,0), B1(0,4), B2(0,8)
  % D_a: OA2PB1
  
  % 塗りつぶし領域 D_a
  \fill[pattern=north east lines, pattern color=gray!60] (0,0) -- (4,0) -- (3,2) -- (0,4) -- cycle;

  % 座標軸
  \draw[->] (-0.5,0) -- (7,0) node[right] {$x$};
  \draw[->] (0,-0.5) -- (0,8.5) node[above] {$y$};
  \node[below left] at (0,0) {$O$};

  % 直線 ell_1: 2x+3y=12 -> 切片 (6,0), (0,4)
  \draw[thick] (6,0) -- (0,4);
  \node[blue, anchor=south west] at (4.5,1.1) {$\ell_1$};
  
  % 直線 ell_2: 2x + y = 8 (a=2) -> 切片 (4,0), (0,8)
  \draw[thick, red] (4,0) -- (0,8);
  \node[red, anchor=south west] at (1,6.5) {$\ell_2$};

  % 定点 P(3,2)
  \fill (3,2) circle (2.5pt);
  \node[above right] at (3,2) {$P$};

  % 頂点ラベル
  \node[below] at (6,0) {$A_1$};
  \node[below, red] at (4,0) {$A_2$};
  \node[left] at (0,4) {$B_1$};
  \node[left, red] at (0,8) {$B_2$};
\end{tikzpicture}
\caption{$\dfrac{4}{3}\le a < \dfrac{8}{3}$の時の$D_a$}
 \end{center}
\end{figure}

この時は$6 \ge \dfrac{8}{a}$だから$A_1$よりも$A_2$が原点に近く，また$\dfrac{8}{4-\frac{3a}{2}} \ge 4$より$B_1$が$B_2$より原点に近いので$D_a$は四角形$OA_2PB_1$になる．
頂点での$g(x,y)$の値は
\begin{align}
 \begin{cases}
  g(A_2) = g\left(\dfrac{8}{a},0\right) = \dfrac{8}{a} \\
  g(P)   = g(3,2) = 5 \\
  g(B_1) = g\left(0,4\right) = 4
 \end{cases}
\end{align}
となるから，最大値は
\begin{align}
 \max g(x,y) 
   &= \max \{g(A_2), g(P)\} \\
   &= 
   \begin{cases}
     5 & \left(\dfrac{8}{5} \le a\right) \\
     \dfrac{8}{a} & \left( a \le \dfrac{8}{5} \right)
   \end{cases} \label{eq:6}
\end{align}
である．

\subparagraph{$3^\circ$ $\dfrac{8}{3}\le a$の時}


\begin{figure}[htb]
\centering

% --- 左図: a = 8/3 の時 ---
\begin{minipage}[b]{0.48\textwidth}
\centering
\begin{tikzpicture}[scale=0.6, >=stealth, every node/.style={font=\scriptsize}]
  % 解析パラメータ (a=8/3) -> A1(6,0), A2(3,0), B1(0,4), B2(存在しない, x=3の垂直線)
  % D_a: OA2PB1 (四角形、ただしA2Pは垂直)
  
  % 塗りつぶし領域 D_a
  \fill[pattern=north east lines, pattern color=gray!60] (0,0) -- (3,0) -- (3,2) -- (0,4) -- cycle;

  % 座標軸
  \draw[->] (-0.5,0) -- (7.5,0) node[right] {$x$};
  \draw[->] (0,-0.5) -- (0,5.5) node[above] {$y$};
  \node[below left] at (0,0) {$O$};

  % 直線 ell_1: 2x+3y=12 -> 切片 (6,0), (0,4)
  \draw[thick, blue] (6,0) -- (0,4);
  \node[blue, anchor=south west] at (4.5,1.1) {$\ell_1$};
  
  % 直線 ell_2: (8/3)x + 0y = 8 -> x = 3 (垂直線)
  \draw[thick, red] (3,-0.5) -- (3,5.0);
  \node[red, anchor=south east] at (3,4.5) {$\ell_2$};

  % 定点 P(3,2)
  \fill (3,2) circle (2.5pt);
  \node[above right] at (3,2) {$P(3,2)$};

  % 頂点ラベル
  \node[below] at (6,0) {$A_1$};
  \node[below left, red] at (3,0) {$A_2(3,0)$};
  \node[left] at (0,4) {$B_1(0,4)$};
  % B2 は存在しない
\end{tikzpicture}
\caption{$a=\dfrac{8}{3}$の時の$D_a$}
\end{minipage}
\hfill
% --- 右図: a > 8/3 の時 (例: a=4) ---
\begin{minipage}[b]{0.48\textwidth}
\centering
\begin{tikzpicture}[scale=0.6, >=stealth, every node/.style={font=\scriptsize}]
  % 解析パラメータ (a=4) -> A1(6,0), A2(2,0), B1(0,4), B2(0,-4) [不採用]
  % D_a: OA2PB1
  
  % 塗りつぶし領域 D_a
  \fill[pattern=north east lines, pattern color=gray!60] (0,0) -- (2,0) -- (3,2) -- (0,4) -- cycle;

  % 座標軸
  \draw[->] (-0.5,0) -- (7.5,0) node[right] {$x$};
  \draw[->] (0,-0.5) -- (0,5.5) node[above] {$y$};
  \node[below left] at (0,0) {$O$};

  % 直線 ell_1: 2x+3y=12 -> 切片 (6,0), (0,4)
  \draw[thick, blue] (6,0) -- (0,4);
  \node[blue, anchor=south west] at (4.5,1.1) {$\ell_1$};
  
  % 直線 ell_2: 4x - 2y = 8 (a=4) -> 切片 (2,0), (0,-4)
  \draw[thick, red] (1.6, -0.8) -- (3.8, 3.6) node[above right] {$\ell_2$};
  \node[red, anchor=south west] at (2.5,2.5) {$\ell_2$};

  % 定点 P(3,2)
  \fill (3,2) circle (2.5pt);
  \node[above right] at (3,2) {$P$};

  % 頂点ラベル
  \node[below] at (6,0) {$A_1$};
  \node[below, red] at (2,0) {$A_2$};
  \node[left] at (0,4) {$B_1$};
 % \node[left, red] at (0,-4) {$B_2$};

\end{tikzpicture}
\caption{$a>\dfrac{8}{3}$の時の$D_a$}
\end{minipage}
\end{figure}

この時は$y$軸上の頂点に関して，$B_2$が存在しないまたは$y<0$の領域にあるため$D_a$の頂点としては$B_1$が採用される．
$x$軸上の点に関しては$6 \ge \dfrac{8}{a}$だから$A_1$よりも$A_2$が原点に近い．
従って$D_a$は四角形$OA_2PB_1$になる．よって$2^{\circ}$の場合と全く同じで\cref{eq:6}より
\begin{align}
 \max g(x,y) = g(P) = 5 \label{eq:7}
\end{align}
である．
\casesend


以上の場合わけで全ての場合が尽くされた．\cref{eq:5,eq:6,eq:7}より
\begin{align}
 \max g(x,y) &=
  \begin{cases}
   g(A_1) = 6 & \left( 0<a\le\dfrac{4}{3} \right) \\
   g(A_2) = \dfrac{8}{a} & \left(\dfrac{4}{3} \le a \le \dfrac{8}{5} \right) \\
   g(P) = 5   & \left(\dfrac{8}{5} \le a\right)
  \end{cases}
\end{align}
が求める最大値である．



\newpage

\end{document}
